\pagebreak
\chapter{Many Mergers Might Fill the Mass Gap \cite{Amen2026} \label{chap:ManyMergersMightFilltheMassGap}}
\hrule
\vspace{.5cm}

\begin{list}{-}{}
\item Stars of mass between 130-250 \(M_\odot\) are thought to end their lives in \gls{ac:PISN}
\item Many BHs in the mass range of 45-120 are still observed
\item The paper proposes that the BHs in this mass range are actually formed via BBH mergers, including \gls{hierarchical-merger}s. 
\item Looking to establish this theory via \gls{ac:GW} observation of hierarchical mergers.
\item 1G + 1G mergers have the same spin axis's and similar masses, as they evolved together. Whereas hierarchical mergers have differing properties
\item They created a model for both non-hierarchical and hierarchical mergers, and then combined said models with a ``mixture fraction'' to create a ``mixture model''\footnote{look into this more and add a full section or something about it :)} Fitting to this model allows for the population to filter itself into its respective subpopulations without enforcing the existence of said subpop. 
\end{list}
\pagebreak

\section{Results}
\hrule
\vspace{.5cm}
\begin{figure}[h!]
    \centering
    \includegraphics[width=.75\textwidth]{ManyMerger.png}
\end{figure}

We see that \dots
\begin{list}{-}{}
\item hierarchical mergers start at \(\sim 46 M_\odot\), supporting that this may fill the mass gap
\item separate collection of hierarchical mergers around 16 solar masses 
\item lack of hierarchical mergers between 20-40 solar masses
\end{list}
